\FloatBarrier
\section{Calibration Feedback}
\label{app:calibration-feedback}

\label{sec:automated-tuning-workflows}

The voltage response in Section~\ref{sec:electron-number-control} and
the pulse settings in Section~\ref{sec:quantum-gates} refer to a chosen
operating point. This supplementary calculation explains how measured
errors can update those settings when the operating point drifts.

A charge transition or spin resonance can drift after a gate pulse has
been calibrated. Correcting one electrode voltage may also shift another
dot or change a tunnel coupling. A feedback update must therefore relate
the measured discrepancies to the physical quantities that each voltage
changes, and its stopping criterion must account for measurement
uncertainty. This appendix develops that update for a two-spin experiment
and for a local vector of device parameters.

The experiment uses the preparation and readout methods of
Appendices~\ref{app:state-preparation} and~\ref{app:spin-readout},
and the spin gates of Appendix~\ref{app:single-spin-control}.

\subsection{Assembling a complete two-spin experiment}

Take two dots, each with a reservoir barrier, and a charge sensor.
As a concrete target, prepare $\ket{01}$ and apply
$\sqrt{\mathrm{SWAP}}$. Under the exchange conditions of
Section~\ref{sec:calibration-two-spin-exchange}, the intended output is
\[
 \ket{01}\longmapsto
 \frac{\ket{01}-i\ket{10}}{\sqrt2}
\]
up to a common phase. This example requires charge control,
initialization, a single-spin flip, exchange, and two spin readouts.
We first establish those controls, then assemble their pulses.

\subsubsection{Establishing the charge settings}
Put the sensor on a current flank and verify its response by loading
and emptying a target dot, using
Section~\ref{sec:measurement-charge-sensor}. With the reservoir
barriers open enough for equilibration, trace target transitions
towards depletion and establish the empty reference before counting
inward to $(1,1)$. A featureless image alone does not identify zero
occupation: a closed barrier can also hide transitions.
The measured transition positions and responses supply the loading
and holding voltages of Appendix~\ref{app:state-preparation}
and the compensated exchange direction of
Appendix~\ref{app:virtual-gates}.

\subsubsection{Obtaining initial spin controls and refining their references}
Load and relax both spins at those settings.
The resonance and Rabi methods of Section~\ref{sec:measurement-rabi}
give the frequency and duration of a spin-flip pulse.
An oscillation period can be estimated with imperfect readout contrast.
The resulting flip then supplies an up-spin reference to improve the
readout voltage and decision rule in Appendix~\ref{app:spin-readout}.
Repeat the spin tests with those improved references.
This is a mutual refinement of preparation, drive and readout, not
an assumption that each was already perfect before the others existed.

Measure $J(v)$ along the compensated path and select the complete
exchange pulse with area $\pi/2$, including the ramps, as in
Section~\ref{sec:calibration-two-spin-exchange}.
These measurements have now supplied all pulse parameters for the
chosen trial.

\subsubsection{Running the trial with fixed settings}
For each repetition, load and relax both spins to $\ket{00}$,
retain the electrons by closing the reservoir barriers, and apply
$R_{x,2}(\pi)$ to prepare $\ket{01}$ up to a common phase.
Apply the compensated exchange pulse, return to weak exchange,
and read both spins with the stored spin-to-charge procedures.
The ideal population prediction is
\[
 P_{01}=P_{10}=\frac12,\qquad P_{00}=P_{11}=0.
\]
For this test, keep the preparation and readout pulses fixed while
varying or correcting the exchange pulse.

Those populations test the transfer but not its relative phase.
The correlation measurements of
Appendix~\ref{app:gate-verification}, with analysis rotations
chosen for the intended output, test its coherence as well.
This separates executing the trial from the extra measurements
needed to verify the gate.

Store the voltage vectors, ramps, hold times, drive frequency and
phase convention, and the readout rule used for this trial.
If a later test fails, remeasure the associated parameter:
a charge-boundary shift changes the voltage setting, a resonance
shift changes the spin-drive reference, and a changed exchange
oscillation changes the exchange pulse. The charge-control stages
can be automated~\cite{Baart2016,Zubchenko2025}; the feedback
calculation below describes how measured errors can drive those
updates.

\subsection{Calculating a voltage correction from measured errors}

Suppose the charge map shows that the first dot's addition energy
is too low, while an exchange measurement indicates insufficient
coupling. Changing one barrier can improve the coupling while
shifting both dot energies. A feedback step uses all measured
responses to make these corrections together.

Near the target, let $\mathbf q(\mathbf V)$ contain an equal
number of independent measured quantities and adjustable voltages.
For example, use the two specified addition-energy offsets relative
to the reservoirs and the interdot tunnel coupling, controlled by
two plungers and one barrier. Their target values encode the
loading or gate condition already chosen above.
Perturb one electrode at a time and remeasure the quantities,
giving the local response
\begin{equation}
 R_{ag}\simeq
 \frac{q_a(\mathbf V+\Delta V_g\mathbf e_g)
       -q_a(\mathbf V-\Delta V_g\mathbf e_g)}
      {2\Delta V_g},
 \label{eq:calibration-feedback-response}
\end{equation}
where $\mathbf e_g$ changes only voltage $g$. Each measurement
uses the charge-transition or tunnel-coupling methods of
Appendices~\ref{app:charge-stability}
and~\ref{app:tunnel-coupling}, respectively.

Write the current discrepancy as
$\mathbf r=\mathbf q^\star-\mathbf q(\mathbf V)$.
For a small change,
\[
 \mathbf q(\mathbf V+\delta\mathbf V)
 =\mathbf q(\mathbf V)+R\delta\mathbf V
                  +O(\|\delta\mathbf V\|^2).
\]
Canceling the first-order discrepancy requires
$R\delta\mathbf V=\mathbf r$. If $R$ is nonsingular, this gives
the feedback step
\begin{equation}
 \mathbf V_{\mathrm{new}}
 =\mathbf V+\eta R^{-1}
          [\mathbf q^\star-\mathbf q(\mathbf V)],
 \qquad 0<\eta\le1.
 \label{eq:calibration-feedback-update}
\end{equation}
The factor $\eta$ limits the step to the measured linear interval
and permitted voltages. Substitution into the expansion gives
\[
 \mathbf r_{\mathrm{new}}
 =(1-\eta)\mathbf r+O(\|\delta\mathbf V\|^2).
\]
Thus $\eta=1$ cancels the first-order residual when the response is
accurate; smaller steps trade a slower correction for a better local
approximation. Remeasure $\mathbf q$ after applying the step.
If the residual does not decrease, reduce the step or determine
$R$ again. A singular $R$ means the selected controls do not give
independent local changes of all the requested quantities.

\subsubsection{A simultaneous energy and tunnel-coupling correction}
As a worked update, take two addition-energy offsets and
$q_3=\ln(t_c/t_{\mathrm{ref}})$, where $t_{\mathrm{ref}}$ is a
fixed energy used only to make the logarithm dimensionless.
Suppose the measured local responses are
\begin{equation}
 \begin{pmatrix}
 \delta\mu_1/(\mu\mathrm{eV})\\
 \delta\mu_2/(\mu\mathrm{eV})\\
 \delta\ln(t_c/t_{\mathrm{ref}})
 \end{pmatrix}
 =
 \begin{pmatrix}
 -100&-25&-10\\
 -20&-80&-15\\
 0.05&0.04&0.50
 \end{pmatrix}
 \begin{pmatrix}
 \delta V_1/(\mathrm{mV})\\
 \delta V_2/(\mathrm{mV})\\
 \delta V_{\mathrm B}/(\mathrm{mV})
 \end{pmatrix}.
 \label{eq:calibration-feedback-example}
\end{equation}
All entries in this displayed, scaled equation are dimensionless.
The first two rows use the addition-energy response from
Appendix~\ref{app:virtual-gates}; the last row
adds illustrative measurements of the tunnel-coupling response.
To request changes $(40,-40,0.10)$ on the left-hand side, solve
the system to obtain
\[
 (\delta V_1,\delta V_2,\delta V_{\mathrm B})
 \simeq(-0.572,\ 0.604,\ 0.209)\,\mathrm{mV}.
\]
This shifts the two energies oppositely while increasing $t_c$
by a factor $e^{0.10}\simeq1.105$. Changing only the barrier
would also shift the dot energies. Apply this full step only if
it stays within the interval used to establish the local response;
otherwise multiply all three changes by the same smaller $\eta$.
Remeasure the two transition positions and the charge-distribution
curve to compute the actual new residual.

\subsubsection{The same update for a pulse angle}
The signed-angle measurement in
Section~\ref{sec:calibration-spin-corrections} supplies a scalar
error instead of a vector of charge and coupling errors.
For a pulse angle, the same first-order calculation gives
\[
 a_{0,\mathrm{new}}=a_0+
 \eta\,\frac{\theta^\star-\theta(a_0)}{d\theta/da_0},
\]
where $a_0$ is the rectangular pulse amplitude, $\theta(a_0)$ comes from
the oscillation measurement, and $\theta^\star$ is the target angle.
Measure the derivative using nearby amplitudes with the same
pulse shape and duration.

\subsection{Deciding whether the correction is resolved and sufficient}

A small measured residual can mean that the control is accurate,
or that the data are too noisy to resolve its error. We therefore
compare the residual with a tolerance tied to the intended operation,
and require the measurement uncertainty to be smaller than that
tolerance.

Before running either loop, specify tolerances $d_a$ for the
measured quantities and require
$|q_a-q_a^\star|\le d_a$ on a new measurement, with measurement
uncertainty small enough to judge that tolerance. Gate-angle or
conditional-phase tolerances are separate from charge-setting
tolerances. If a later check fails, repeat the measurement associated
with that failure: a shifted charge boundary requires retuning the
voltage setting; a shifted spin resonance requires updating the
drive reference. Changing a setting also invalidates downstream
calibrations whose Hamiltonian parameters changed. This dependency,
rather than a universal retuning schedule, determines what must be
measured again.

The tolerance should be connected to the operation being requested.
For a coherent angle error $\eta_\theta$ about the intended axis,
Eq.~\eqref{eq:calibration-average-fidelity} gives
$1-F_{\mathrm{avg}}\simeq\eta_\theta^2/6$. Allocating an error
$r_{\mathrm{angle}}$ to this mechanism therefore sets the
angle tolerance
$|\eta_\theta|\lesssim\sqrt{6r_{\mathrm{angle}}}$.
In the signed sequence of
Eq.~\eqref{eq:calibration-signed-angle-error}, $n$ trials near
$P_1=1/2$ have counting uncertainty approximately
$1/(2\sqrt n)$. The inferred angle uncertainty is then
$1/(m\sqrt n)$ before assignment corrections and other errors.
This relation determines how many repetitions can resolve the
chosen tolerance. Increasing $m$ helps only while the repeated
sequence still obeys the small-error model and retains contrast.
The stopping test consequently concerns both the measured residual
and whether the measurement is precise enough to resolve it.
